\documentclass[11pt]{article}
\usepackage{amsmath,amssymb,amsthm}
\newtheorem{theorem}{Theorem}
\title{Exterior Query Compactification and Exact Moment Remainders}
\author{}
\date{November 24, 2026}
\begin{document}
\maketitle

\paragraph{Status.}
The identities and the two cube-admissibility statements below are
Lean-checked in \texttt{ReverseCauchyCompactification.lean} and
\texttt{ReverseMomentGerms.lean}, in namespace
\texttt{HeadComplexity.AxisZeroGluing}. They use the previously checked
paired-axis construction in \texttt{ReverseMaskedCauchy.lean}.
They do not assert a universal sign decoder. Neither accepted
asymptotic target follows without a further theorem.

Let $p_{hi}>0$ be row-dependent poles, $a_{hi}$ arbitrary residues,
$t_j>p_{hi}$ positive query times, and $M_{hj}\in\{0,1\}$ masks.
For $\varepsilon>0$, put
\[
 \phi(x)=\frac{x}{x+\varepsilon},\qquad
 p'_{hi}=\phi(p_{hi}),\quad t'_j=\phi(t_j),\quad
 a'_{hi}=\frac{\varepsilon a_{hi}}{p_{hi}+\varepsilon}.
\]
The exact scalar identities are
\[
 \phi(t)-\phi(p)
 =\frac{\varepsilon(t-p)}{(t+\varepsilon)(p+\varepsilon)},
 \qquad
 (1-\phi(t))\frac{\varepsilon a/(p+\varepsilon)}{\phi(t)-\phi(p)}
 =\varepsilon\frac a{t-p}.
 \tag{1}
\]
The second identity includes the factor from the paired-axis
construction. Consequently the net score scaling is independent of
the query, not merely positive at each query.

\begin{theorem}[Exact compactified axis realization]
Suppose, for every head $h$,
\[
 \sum_i\frac{p_{hi}}{p_{hi}+\varepsilon}
 +\sum_j M_{hj}\frac{\varepsilon}{t_j+\varepsilon}<1.
 \tag{2}
\]
There are $2H$ affine-ratio heads with strictly negative denominator
slopes and positive denominators on the entire Boolean $(r+s)$-cube.
Their sum is exactly zero on both Boolean axes and has grid values
\[
 F(e_i,e_j)=\varepsilon\sum_h M_{hj}\frac{a_{hi}}{t_j-p_{hi}}.
 \tag{3}
\]
\end{theorem}
\begin{proof}
All transformed poles satisfy $0<p'_{hi}<t'_j<1$.
The strict inequality (2) persists, for sufficiently small $\eta>0$,
as
\[
 \sum_i p'_{hi}+s\eta+\sum_jM_{hj}(1-t'_j)<1+\eta.
\]
There are only finitely many heads, so a common positive $\eta$ works.
Apply the checked paired-axis construction with denominator scale
one, strict padding $\eta$, poles $p'$, residues $a'$, and times $t'$.
Its value is
$(1-t'_j)\sum_h M_{hj}a'_{hi}/(t'_j-p'_{hi})$.
Identity (1), term by term, proves (3). All axis and full-cube
conditions are supplied by that construction.
\end{proof}

The Lean declaration is
\texttt{exists\_compactified\_maskedCauchy\_reverseAxisData}.

\begin{theorem}[A sufficient small-pole budget]
Let $C,b,\varepsilon>0$. It is enough to have
\[
 0<p_{hi}\leq\varepsilon^2 C<b\leq t_j,
 \qquad \varepsilon(rC+s/b)<1.
 \tag{4}
\]
\end{theorem}
\begin{proof}
One has $p/(p+\varepsilon)\leq\varepsilon C$ and
$\varepsilon/(t+\varepsilon)\leq\varepsilon/b$.
Summing proves (2). The strict separation of poles and queries also
follows from (4).
\end{proof}
The declaration is
\texttt{exists\_compactified\_maskedCauchy\_of\_small\_poles}.
In particular, fixed positive query centers may be separated from
one another. They need not cluster near one. This observation does
not permit poles to approach interior query points: the budget
(2), or a sufficient condition such as (4), must still be proved.

\paragraph{Exact moments.}
For $t\neq0$, $t-\delta p_h\neq0$, and vanishing moments
$\sum_h a_hp_h^\ell=0$ for every $\ell<k$, the exact formula is
\[
 \sum_h\frac{a_h}{t-\delta p_h}
 =\left(\frac\delta t\right)^k
   \sum_h\frac{a_hp_h^k}{t-\delta p_h}.
 \tag{5}
\]
The proof subtracts one geometric term and inducts on $k$.
It is checked as \texttt{cauchy\_moment\_remainder}; its one-step
form is \texttt{cauchy\_moment\_step}.

For the cubic family used in the four-head, ten-query written
construction, set $D_h=t-\delta P_h-\delta^2P_{2h}$.
If $tD_h\neq0$, $\sum_hA_h=0$, and $\sum_hA_hP_h=0$, then
\[
 \sum_h\frac{\delta A_h+\delta^3 A_{3h}}{D_h}
 =\delta^3\left(
 \sum_h\frac{A_{3h}}{D_h}
 +\frac{\sum_hA_hP_{2h}}{t^2}
 +\sum_h\frac{A_h(P_h+\delta P_{2h})^2}{t^2D_h}
 \right).
 \tag{6}
\]
This is \texttt{cauchy\_cubic\_remainder}. Its proof is an exact
per-head rational identity followed by the two cancellations.
The ten-query sign certificates and their universality argument
remain written and computational results, not a consequence claimed
solely from (6).

\paragraph{Verification boundary.}
The aggregate \texttt{HeadComplexity.Results.All} builds successfully
with these modules. No original model definition was changed.
The central missing construction is still a sign-covering family
whose number of queries per head tends to infinity at the required
rate, or a balanced joint-table compiler with the corresponding
saving. The compactification theorem reduces an admissibility issue;
it does not prove that sign-covering statement.
\end{document}
